\section{Algebralliset käyrät ja projektiiviset avaruudet}
Tämän osion tarkoituksena on esitellä tiiviisti algebrallisten käyrien, projektiivisten avaruuksien ja projektiivisten käyrien määritelmän sanaston kiinnittämiseksi. Kiinnostunut lukija löytää huolellisemman tarkastelun esimerkiksi Gathmannin algebrallisen geometrian luentomuistiinpanoista \citep{gathmann-ag}.

\subsection{Affiinit algebralliset käyrät kompleksitasossa}
Affiini algebrallinen käyrä kompleksitasossa on määritelty polynomin nollakohtien joukkona. Tarkemmin sanottuna, jos $f \in \mathbb{C}[x,y]$ on ei-vakio polynomi, niin affiini algebrallinen käyrä $C$ joka liittyy polynomiin $f$ on joukko
\begin{equation}
    C = \{ (x,y) \in \mathbb{C}^2 \mid f(x,y) = 0 \}.
\end{equation}
Koska $\C$ on algebrallisesti suljettu, niin jokainen polynomi $f \in \mathbb{C}[x,y]$ voidaan tekijöidä lineaarisiin tekijöihin. Tämän vuoksi jokainen affiini algebrallinen käyrä kompleksitasossa voidaan esittää lineaaristen käyrien yhdistelmänä.

\subsection{Singulariteetit ja säännöllisyys}
Käyrän $C$ paikallisen muodon pisteessä $p=(a,b)\in C$ määrittyy $f$:n osittaisderivaattojen avulla. Piste $p$ on \emph{singulaarinen} (tai kriittinen), jos $f$:n osittaisderivaatat häviävät pisteessä $p$:
\begin{equation}
    \frac{\partial f}{\partial x}(a,b) = 0 \quad \text{ja} \quad \frac{\partial f}{\partial y}(a,b) = 0
\end{equation}
Jos vähintään yksi osittaisderivaatoista ei ole nolla $p$:ssä, niin piste on \emph{säännöllinen}. Säännöllisessä pisteessä implisiittisten funktioiden lause takaa, että käyrä näyttää paikallisesti kompleksisen 1-ulotteiselta varistolta, ja on mahdollista määritellä yksikäsitteinen tangenttisuora.

\subsection{Projektiivinen taso ja suora}
Affiinilla avaruudella $\C$ on geometrinen rajoitus: kaksi erillistä suoraa leikkaavat yleensä yhdessä pisteessä, paitsi jos ne ovat yhdensuuntaisia. Yhtenäistääksemme näitä geometrisia käsitteitä, otamme käyttöön projektiivisen avaruuden.

Kompleksinen $n$-ulotteinen projektiivinen avaruus $\P^n(\C)$ määritelläänjoukoksi, joka koostuu avaruuden $\C^{n+1}$ vektorisäteistä (tai niiden ekvivalenssiluokista), jotka kulkevat origon kautta. Toisin sanoen, $\P^n(\C)$ on määritelty seuraavasti:
\begin{maaritelma}
    Olkoon $(Z_0,\dots,Z_n) \in \C^{n+1} \setminus \{0\}$. Määritellään ekvivalenssirelaatio $\sim$ seuraavasti:
    \begin{equation*}
        (Z_0,\dots,Z_n) \sim (\lambda Z_0, \dots, \lambda Z_n) \quad \text{kaikille } \lambda \in \C^*.
    \end{equation*}
    Tällöin
    \begin{equation*}
        \P^n(\C):=\frac{\C^{n+1}\setminus\{0\}}{\sim}.
    \end{equation*}
    Ekvivalenssiluokan edustaja merkitään homogeenisilla koordinaateilla $[Z_0:\dots:Z_n]$.
\end{maaritelma}
\begin{asetus}
    Jatkossa merkitään $\P^n$ pidemmän $\P^n(\C)$ sijasta, koska tässä työssä keskitytään yksinomaan kompleksisten käyrien tarkasteluun.
\end{asetus}
\begin{esimerkki}
    \label{esim:proj-taso-ja-suora}
    Tapauksissa $n=1$ ja $n=2$
    \begin{itemize}
        \item $\P^1\cong\C\cup\{\infty\}$ (Riemann-pallo)
        \item $\P^2\cong\C^2\cup\P^1$. $\P^2$:n koordinaatit ovat $[X:Y:Z]$, ja $\P^1$ on joukko pisteitä, joilla $Z=0$.
    \end{itemize}
    Jos tarkastellaan projektiivisten avaruuksien ulottuvuuksia reaalimielessä ($\dim_\R(\C)=2$), niin $\dim_\R(\P^1)=2$ ja $\dim_\R(\P^2)=6$ (vai 4?).
\end{esimerkki}

\subsection{Upotus \texorpdfstring{$\C^2\hookrightarrow\P^2$}{Upotus C²→P²}}
\begin{lemma}
    \label{lemma:projektiivisen-avaruuden-jako}
    Projektiivinen avaruus on erillinen yhdiste $\P^{n}=\C^n\sqcup\P^{n-1}$.
\end{lemma}
\begin{proof}
    Tarkastellaan $\P^{n}$:n pistettä $[Z_0:\dots:Z_{n}]$, jolla on $n$ koordinaattia. Jos $Z_{n+1}\neq 0$, niin piste voidaan esittää muodossa
    \begin{equation*}
        [Z_0:\dots:Z_n] \sim [Z_0:\dots:Z_{n-1}:1] \in \C^n,
    \end{equation*}
    koska on olemassa $\lambda\in\C\setminus\{0\}$, jolla $\lambda Z_n=1$. Jos taas $Z_n=0$, niin piste kuuluu $\P^{n-1}$-joukkoon. Tällöin saadaan jako $\P^{n}=\C^n\sqcup\P^{n-1}$.
\end{proof}
Lemma yleistää esimerkin \ref{esim:proj-taso-ja-suora}: yleisen projektiivisen avaruuden $\P^n$ ``sisus'' on $\C^n$ ja ``kuori äärettömyydessä'' $\P^{n-1}$, joka induktiivisesti kuortaan lukuunottamatta muistuttaa $\C^{n-1}$:tä. Tämä antaa mielekkään upotuksen $\C^n\hookrightarrow\P^n$.

Työlle relevantissa 2-ulotteisessa tapauksessa avaruus $\C^2$ voidaan nähdä paikallisena karttana (tai tiheänä avoimena joukkona) projektiivisessa tasossa $\P^2(\C)$. Tyypillinen upotus määritellään seuraavasti:
\begin{equation}
    \begin{aligned}
        \iota : \C^2 &\hookrightarrow \P^2(\C) \\
        (x,y) &\mapsto [x : y : 1]
    \end{aligned}
\end{equation}
Pisteet $\P^2$:ssa, jotka eivät kuulu tämän upotuksen kuvaan, ovat muotoa $[X:Y:0]$. Näiden pisteiden joukko muodostaa projektiivisen suoran $\P^1$, jota kutsutaan \emph{suoraksi äärettömyydessä} tai \emph{äärettömyyssuoraksi}. Lemman \ref{lemma:projektiivisen-avaruuden-jako} hengessä $\P^2=\C^2\sqcup\P^1$.

\subsection{Projektiiviset käyrät}
Projektiivisessa tasossa $\P^2$ klassinen polynomi $f(x,y)$ ei ole hyvin määritelty, koska homogeeniset koordinaatit ovat ekvivalentteja skalaarikertoimen suhteen. Tämän vuoksi on käytettävä \emph{homogeenistä} polynomia $F(X,Y,Z)$, eli polynomia, jossa kaikilla monomeilla on sama aste $d$. Yhtälö $F(\lambda X, \lambda Y, \lambda Z) = \lambda^d F(X,Y,Z) = 0$ on tällöin riippumaton valitusta edustajasta.
\begin{maaritelma}[Homogenisointi]
    Polynomi $f(x,y) \in \C[x,y]$ voidaan \emph{homogenisoida} määrittelemällä
    \begin{equation}
        F(X,Y,Z) = Z^d f\left(\frac{X}{Z}, \frac{Y}{Z}\right),
    \end{equation}
    missä $d$ on polynomin $f$ aste.
\end{maaritelma}
\begin{maaritelma}
    Projektiivinen algebrallinen käyrä $\mathcal{C}$ määritellään homogeenisen polynomin $F(X,Y,Z)$ nollakohtien joukkona:
    \begin{equation}
        \mathcal{C} = \{ [X:Y:Z] \in \P^2 \mid F(X,Y,Z) = 0 \}.
    \end{equation}
\end{maaritelma}
\begin{huomautus}
    Polynomin homogenisointi sallii luonnollisesti käyrän tarkastelun myös äärettömyyssuoralla. Projektiivinen käyrä $\mathcal{C}$ on aina suljettu ja kompakti, koska se on suljettu alijoukko projektivisessa avaruudessa $\P^2$, joka on itsessään kompakti.
\end{huomautus}
